% GENERATED FILE. Edit canonical lessons, then run npm run generate:books.
% canonical-source-hash: fnv1a64:a2d6b153afd3acd3
% canonical-lessons: GE-C359-gerechtigkeit, GE-C359-ungleichheit, GE-C359-armut, GE-C359-reichtum, GE-C359-benachteiligen

\chapter{Gerecht oder nicht?}
\label{ch:ge-gerecht-oder-nicht}
\hlchaptermodality{359}

\begin{chapteropening}
\textbf{By the end of this chapter:} \emph{I can talk about justice and inequality, poverty and wealth, and people who are put at a disadvantage.}

Complete the last lesson of chapter 359: I can talk about justice and inequality, poverty and wealth, and people who are put at a disadvantage.
\end{chapteropening}

\section[die Gerechtigkeit]{die Gerechtigkeit — justice, fairness}

\label{lesson:GE-C359-gerechtigkeit}



\begin{quote}
Before the new one: say the German for the vote, then the German for to contribute.
\end{quote}

\subsection*{You'll want to know: die Gerechtigkeit}
\textbf{die Gerechtigkeit} — \textquotedblleft{}justice, fairness\textquotedblright{}.

\begin{quote}
Viele junge Leute kämpfen für mehr Gerechtigkeit.
\end{quote}

\begin{quote}
\emph{Many young people fight for more justice.}
\end{quote}

No plural in everyday use. From \emph{gerecht} — fair.

\begin{cousinweb}
\emph{Gerechtigkeit} is formed from \emph{gerecht} \textquotedblleft{}fair, righteous\textquotedblright{}, Old High German \emph{gireht} (EtymDB), with \emph{-igkeit}; the adjective is built on the word for \textquotedblleft{}right\textquotedblright{}, Proto-Germanic \emph{rehtaz}.
\end{cousinweb}

\begin{grammarlens}[title={What you've built}]
One more word for telling, explaining and describing at B1.
\end{grammarlens}

\subsection*{Your turn}
\begin{itemize}
\raggedright
  \item \emph{Say it:} \emph{die Gerechtigkeit}
  \item \emph{Say it:} \emph{die Gerechtigkeit}, once more
  \item \emph{Say it:} a sentence of your own with \emph{die Gerechtigkeit}
  \item \emph{Recall:} say the German for the vote, then the German for to contribute, then say \emph{die Gerechtigkeit} again
\end{itemize}

\begin{tcolorbox}[breakable,colback=teal!4,colframe=teal!35!black,title={Before you move on}]
What does die Gerechtigkeit mean? (\textquotedblleft{}Justice, fairness\textquotedblright{}.) Say it once more.
\end{tcolorbox}
\section[die Ungleichheit]{die Ungleichheit — inequality}

\label{lesson:GE-C359-ungleichheit}



\begin{quote}
Before the new one: say the German for to contribute, then the German for justice.
\end{quote}

\subsection*{You'll want to know: die Ungleichheit}
\textbf{die Ungleichheit} — \textquotedblleft{}inequality\textquotedblright{}.

\begin{quote}
Die Ungleichheit zwischen Arm und Reich wird größer.
\end{quote}

\begin{quote}
\emph{The inequality between rich and poor is growing.}
\end{quote}

Plural: \emph{die Ungleichheiten}. From \emph{ungleich} — unequal.

\begin{cousinweb}
\emph{Ungleichheit} is built from \emph{un-} \textquotedblleft{}not\textquotedblright{}, \emph{gleich} \textquotedblleft{}equal\textquotedblright{}, inherited from Old High German \emph{gilīh} and Proto-Germanic \emph{galīkaz} \textquotedblleft{}alike\textquotedblright{} (EtymDB), and the ending \emph{-heit}.
\end{cousinweb}

\begin{grammarlens}[title={What you've built}]
One more word for telling, explaining and describing at B1.
\end{grammarlens}

\subsection*{Your turn}
\begin{itemize}
\raggedright
  \item \emph{Say it:} \emph{die Ungleichheit}
  \item \emph{Say it:} \emph{die Ungleichheit}, once more
  \item \emph{Say it:} a sentence of your own with \emph{die Ungleichheit}
  \item \emph{Recall:} say the German for to contribute, then the German for justice, then say \emph{die Ungleichheit} again
\end{itemize}

\begin{tcolorbox}[breakable,colback=teal!4,colframe=teal!35!black,title={Before you move on}]
What does die Ungleichheit mean? (\textquotedblleft{}Inequality\textquotedblright{}.) Say it once more.
\end{tcolorbox}
\section[die Armut]{die Armut — poverty}

\label{lesson:GE-C359-armut}



\begin{quote}
Before the new one: say the German for justice, then the German for inequality.
\end{quote}

\subsection*{You'll want to know: die Armut}
\textbf{die Armut} — \textquotedblleft{}poverty\textquotedblright{}.

\begin{quote}
In vielen Ländern leben Kinder in großer Armut.
\end{quote}

\begin{quote}
\emph{In many countries children live in great poverty.}
\end{quote}

No plural. Built on \emph{arm} — poor.

\begin{cousinweb}
\emph{Armut} is inherited from Middle High German \emph{armuot} (EtymDB), built on \emph{arm} \textquotedblleft{}poor\textquotedblright{} with an old noun ending; EtymDB lists Dutch \emph{armoede} as a cousin.
\end{cousinweb}

\begin{grammarlens}[title={What you've built}]
One more word for telling, explaining and describing at B1.
\end{grammarlens}

\subsection*{Your turn}
\begin{itemize}
\raggedright
  \item \emph{Say it:} \emph{die Armut}
  \item \emph{Say it:} \emph{die Armut}, once more
  \item \emph{Say it:} a sentence of your own with \emph{die Armut}
  \item \emph{Recall:} say the German for justice, then the German for inequality, then say \emph{die Armut} again
\end{itemize}

\begin{tcolorbox}[breakable,colback=teal!4,colframe=teal!35!black,title={Before you move on}]
What does die Armut mean? (\textquotedblleft{}Poverty\textquotedblright{}.) Say it once more.
\end{tcolorbox}
\section[der Reichtum]{der Reichtum — wealth, riches}

\label{lesson:GE-C359-reichtum}



\begin{quote}
Before the new one: say the German for inequality, then the German for poverty.
\end{quote}

\subsection*{You'll want to know: der Reichtum}
\textbf{der Reichtum} — \textquotedblleft{}wealth, riches\textquotedblright{}.

\begin{quote}
Sein Reichtum hat ihn nicht glücklich gemacht.
\end{quote}

\begin{quote}
\emph{His wealth did not make him happy.}
\end{quote}

Plural: \emph{die Reichtümer}. From \emph{reich} — rich.

\begin{cousinweb}
\emph{Reichtum} is inherited from Middle High German \emph{rīchtuom} and Old High German \emph{rīhhituom} (EtymDB), built on \emph{reich} \textquotedblleft{}rich\textquotedblright{} with \emph{-tum}; EtymDB lists Dutch \emph{rijkdom} and English \emph{richdom} as cousins.
\end{cousinweb}

\begin{grammarlens}[title={What you've built}]
One more word for telling, explaining and describing at B1.
\end{grammarlens}

\subsection*{Your turn}
\begin{itemize}
\raggedright
  \item \emph{Say it:} \emph{der Reichtum}
  \item \emph{Say it:} \emph{der Reichtum}, once more
  \item \emph{Say it:} a sentence of your own with \emph{der Reichtum}
  \item \emph{Recall:} say the German for inequality, then the German for poverty, then say \emph{der Reichtum} again
\end{itemize}

\begin{tcolorbox}[breakable,colback=teal!4,colframe=teal!35!black,title={Before you move on}]
What does der Reichtum mean? (\textquotedblleft{}Wealth, riches\textquotedblright{}.) Say it once more.
\end{tcolorbox}
\section[benachteiligen]{benachteiligen — to put at a disadvantage}

\label{lesson:GE-C359-benachteiligen}



\begin{quote}
Before the new one: say the German for poverty, then the German for wealth.
\end{quote}

\subsection*{You'll want to know: benachteiligen}
\textbf{benachteiligen} — \textquotedblleft{}to put at a disadvantage\textquotedblright{}.

\begin{quote}
Das neue Gesetz benachteiligt Familien mit wenig Geld.
\end{quote}

\begin{quote}
\emph{The new law puts families with little money at a disadvantage.}
\end{quote}

Inseparable: \emph{es benachteiligt, es hat benachteiligt}. \emph{Benachteiligt sein} — to be at a disadvantage.

\begin{cousinweb}
\emph{benachteiligen} is built with \emph{be-} on \emph{Nachteil} \textquotedblleft{}disadvantage\textquotedblright{}, from \emph{nach} \textquotedblleft{}after\textquotedblright{} and \emph{Teil} \textquotedblleft{}part\textquotedblright{}, Proto-Germanic \emph{dailą} (EtymDB): you leave someone the worse part.
\end{cousinweb}

\begin{grammarlens}[title={What you've built}]
One more word for telling, explaining and describing at B1.
\end{grammarlens}

\subsection*{Your turn}
\begin{itemize}
\raggedright
  \item \emph{Say it:} \emph{benachteiligen}
  \item \emph{Say it:} \emph{benachteiligen}, once more
  \item \emph{Say it:} a sentence of your own with \emph{benachteiligen}
  \item \emph{Recall:} say the German for poverty, then the German for wealth, then say \emph{benachteiligen} again
\end{itemize}

\begin{tcolorbox}[breakable,colback=teal!4,colframe=teal!35!black,title={Before you move on}]
What does benachteiligen mean? (\textquotedblleft{}To put at a disadvantage\textquotedblright{}.) Say it once more.
\end{tcolorbox}
